\documentclass[letterpaper]{article} % DO NOT CHANGE THIS
\usepackage{aaai24}  % DO NOT CHANGE THIS
\usepackage{times}  % DO NOT CHANGE THIS
\usepackage{helvet}  % DO NOT CHANGE THIS
\usepackage{courier}  % DO NOT CHANGE THIS
\usepackage[hyphens]{url}  % DO NOT CHANGE THIS
\usepackage{graphicx} % DO NOT CHANGE THIS
\urlstyle{rm} % DO NOT CHANGE THIS
\def\UrlFont{\rm}  % DO NOT CHANGE THIS
\usepackage{natbib}  % DO NOT CHANGE THIS AND DO NOT ADD ANY OPTIONS TO IT
\usepackage{caption} % DO NOT CHANGE THIS AND DO NOT ADD ANY OPTIONS TO IT
\frenchspacing  % DO NOT CHANGE THIS
\setlength{\pdfpagewidth}{8.5in}  % DO NOT CHANGE THIS
\setlength{\pdfpageheight}{11in}  % DO NOT CHANGE THIS
\usepackage{amsmath}
\usepackage{placeins}  % Fix the random figure placement by LaTeX
\usepackage{lipsum}
\usepackage{enumitem}

% DISALLOWED PACKAGES
% \usepackage{authblk} -- This package is specifically forbidden
% \usepackage{balance} -- This package is specifically forbidden
% \usepackage{color (if used in text)
% \usepackage{CJK} -- This package is specifically forbidden
% \usepackage{float} -- This package is specifically forbidden
% \usepackage{flushend} -- This package is specifically forbidden
% \usepackage{fontenc} -- This package is specifically forbidden
% \usepackage{fullpage} -- This package is specifically forbidden
% \usepackage{geometry} -- This package is specifically forbidden
% \usepackage{grffile} -- This package is specifically forbidden
% \usepackage{hyperref} -- This package is specifically forbidden
% \usepackage{navigator} -- This package is specifically forbidden
% (or any other package that embeds links such as navigator or hyperref)
% \indentfirst} -- This package is specifically forbidden
% \layout} -- This package is specifically forbidden
% \multicol} -- This package is specifically forbidden
% \nameref} -- This package is specifically forbidden
% \usepackage{savetrees} -- This package is specifically forbidden
% \usepackage{setspace} -- This package is specifically forbidden
% \usepackage{stfloats} -- This package is specifically forbidden
% \usepackage{tabu} -- This package is specifically forbidden
% \usepackage{titlesec} -- This package is specifically forbidden
% \usepackage{tocbibind} -- This package is specifically forbidden
% \usepackage{ulem} -- This package is specifically forbidden
% \usepackage{wrapfig} -- This package is specifically forbidden

% DISALLOWED COMMANDS
% \nocopyright -- Your paper will not be published if you use this command
% \addtolength -- This command may not be used
% \balance -- This command may not be used
% \baselinestretch -- Your paper will not be published if you use this command
% \clearpage -- No page breaks of any kind may be used for the final version of your paper
% \columnsep -- This command may not be used
% \newpage -- No page breaks of any kind may be used for the final version of your paper
% \pagebreak -- No page breaks of any kind may be used for the final version of your paperr
% \pagestyle -- This command may not be used
% \tiny -- This is not an acceptable font size.
% \vspace{- -- No negative value may be used in proximity of a caption, figure, table, section, subsection, subsubsection, or reference
% \vskip{- -- No negative value may be used to alter spacing above or below a caption, figure, table, section, subsection, subsubsection, or reference

% Keep the \pdfinfo as shown here. No need to add /Title or /Author tags.
\pdfinfo{
/TemplateVersion (2024.1)
}

\setcounter{secnumdepth}{1}

% Title should use Title Case.
\title{\huge AI6103 - Deep Learning Group Project Report \\ \Large Revisiting Sentence-BERT: Reproduction and Lightweight Enhancements}

\author{
\begin{tabular}{ccc}
\textbf{Dai Zi Qiao} & \textbf{Dai Dongjie} & \textbf{Li Huilin} \\
G2509604K & G2503624K& G2503458K \\
ZDAI002@e.ntu.edu.sg & DONGJIE001@e.ntu.edu.sg & HUILIN002@e.ntu.edu.sg \\[1em]
\multicolumn{3}{c}{
\begin{tabular}{cc}
\textbf{Li Zian} & \textbf{Zhang Yanhai} \\
G2509698B & G2509626L \\
C250276@e.ntu.edu.sg & ZH0010AI@e.ntu.edu.sg \\
\end{tabular}
}
\end{tabular}
}

\affiliations{
    \large Nanyang Technological University \\[1em]
}

\begin{document}

\maketitle

\begin{abstract}
This project revisits the original Sentence-BERT framework by Reimers and Gurevych \cite{reimers-gurevych-2019-sentence-bert} and evaluates its reproducibility in a legacy implementation setting. 
Our goal is to reproduce the original SBERT-base training pipeline and investigate whether lightweight objective-level modifications can improve Semantic Textual Similarity (STS) performance on the STSBenchmark dataset without changing the underlying architecture. 
The training pipeline first warm-starts the encoder on natural language inference (NLI), then optionally reshapes the embedding space with an intermediate objective, and finally adapts the model to STSBenchmark with direct similarity supervision. 
Based on this pipeline, we study two main enhancement directions: AoE-style STS fine-tuning, which sharpens local similarity ranking, and an MNRL-based contrastive stage with a lightweight uniformity regularizer, which improves the global embedding distribution. Furthermore, we conducted a comparative study between standard Mean Pooling and an Attention Pooling mechanism to evaluate if structural complexity could offer further gains.
According to the experimental results, the best result is achieved by applying better contrastive learning with uniformity regularizer, which obtains the highest cosine Spearman score on the dataset. 
Overall, the results show that within the 2019 SBERT codebase, meaningful gains can still be achieved through careful objective design, runtime configuration, and evaluation without modifying the original architecture.

\end{abstract}

\section{Introduction}
Sentence-BERT was proposed to address the inefficiency of standard BERT on sentence-pair tasks. Instead of jointly encoding every sentence pair at inference time, it extends BERT with a siamese architecture that produces a fixed-size sentence embedding for each input. 
These embeddings can then be compared efficiently using cosine similarity. 
In this way, SBERT becomes better suited to Semantic Textual Similarity (STS) and related retrieval-style applications.

Although sentence embedding models have advanced substantially, many practical workflows still depend on older training stacks because of compute constraints or compatibility requirements. This motivates a narrower but important question: can we still obtain meaningful gains within the original SBERT framework without introducing architectural change?

To answer this question, we keep the SBERT backbone unchanged and investigate objective-level modifications under a three-stage curriculum: NLI warm-start, optional geometry shaping, and final STS adaptation. The plain reference systems in this report are \texttt{STS} and \texttt{NLI+STS}. 
On top of them, we study two lightweight enhancement directions: an AoE-style ranking objective for finer local similarity ordering, and an MNRL-based contrastive stage with UniReg for improving the global geometry of the embedding space.

The final results show that objective design matters more than architectural modification in this legacy setting. All three main enhanced branches outperform the plain baselines, with \texttt{NLI+MNRL+UniReg+STS} achieving the highest mean cosine Spearman score on STSBenchmark. 
At the same time, some seemingly reasonable combinations, such as jointly applying MNRL and AoE, do not yield further gains. These findings suggest that even within a 2019-style SBERT codebase, careful objective design can still produce meaningful and reproducible improvements.

\section{Related Work}

\subsection{The Foundation of Efficient Sentence Embeddings}
Before the widespread adoption of transformer-based siamese networks, models like the Universal Sentence Encoder \cite{cer-etal-2018-universal} provided early, robust baselines for encoding sentences into fixed-length vectors. 
However, our research builds directly upon Sentence-BERT (SBERT), which adapted the standard BERT architecture into a siamese and triplet network framework \cite{reimers-gurevych-2019-sentence-bert, devlin-etal-2019-bert}. 
While the field has since explored decoder-only architectures for text embeddings, such as SGPT \cite{muennighoff-2022-sgpt}, our implementation strictly preserves the 2019-style SBERT training stack. 
Rather than introducing architectural shifts, we utilize a three-stage sentence embedding curriculum---warm-starting on Natural Language Inference (NLI) data, applying an intermediate geometry-shaping objective, and adapting to Semantic Textual Similarity (STS) data---to achieve lightweight, reproducible gains.

\subsection{The Shift Towards Contrastive Representation Learning}
To improve the global geometry of the embedding space, our methodology intersects with the modern paradigm of contrastive learning. 
The introduction of SimCSE \cite{gao-etal-2021-simcse} marked a significant shift by demonstrating how standard supervised contrastive learning could drastically improve embedding alignment. 
Subsequent frameworks built heavily upon this foundation: DiffCSE \cite{chuang-etal-2022-diffcse} introduced equivariant contrastive learning to capture subtle semantic differences. 
Meanwhile, modern state-of-the-art models like E5 \cite{wang-etal-2022-e5}, GTE \cite{li-etal-2023-gte}, and M3 Embedding \cite{chen2024bge} demonstrated the sheer power of multi-stage contrastive pre-training and self-knowledge distillation over massive, multi-lingual text pairs. 

Inspired by these multi-stage curricula, our approach utilizes Multiple Negatives Ranking Loss (MNRL), which functions as a sentence-pair instantiation of the InfoNCE objective \cite{rusak2025infonce}. 
To counteract dimensional collapse, we pair this with a lightweight uniformity regularizer (UniReg) inspired by contrastive alignment properties \cite{wang-isola-2020-alignment-uniformity}. 
This encourages embeddings to spread evenly across the unit hypersphere without requiring the massive parameter scale of modern alternatives.

\subsection{Angle-Optimized Similarity Supervision}
Standard cosine regression frequently suffers from a saturation zone, where gradient signals diminish for highly similar sentence pairs. To address this limitation in the final STS fine-tuning stage, we draw upon Angle-optimized Embeddings (AoE) \cite{li-li-2024-aoe}. 
Instead of strictly relying on saturated cosine regression, AoE introduces a ranking-aware supervision signal based on angular relations. By replacing the plain regression objective with an angle-based ordering on graded pairs, our model achieves heightened sensitivity to fine-grained semantic differences without requiring structural modifications to the underlying BERT encoder.

\section{Experiment Setups}
We evaluate all methods within a single reproducible pipeline built on the legacy
\texttt{sentence-transformers v0.2.4} stack. The backbone is fixed to
\texttt{bert-base-uncased} with mean pooling, so the main experimental variables are the training
curriculum, intermediate objective, and final STS supervision.

\subsection{Execution Environment}
Experiments were run on GPU servers under a unified software environment with fixed package
versions. This standardization is important because the project was executed across multiple
machines, and we wanted observed differences to reflect training objectives rather than environment
drift.

\subsection{Model Configurations}
We compare two plain baselines and several lightweight enhancement branches:
\begin{itemize}[itemsep=0pt]
    \item \textbf{Plain baselines:} \texttt{STS} and \texttt{NLI+STS}.
    \item \textbf{Contrastive branch:} \texttt{NLI+MNRL+STS}, with \texttt{NLI+MNRL+UniReg+STS} as the regularized variant.
    \item \textbf{Ranking branch:} \texttt{NLI+AoE+STS}, which replaces plain STS cosine regression with AoE-style ranking supervision.
    \item \textbf{Additional variants:} joint \texttt{MNRL+AoE}, CoSENT-based STS replacement, and attention pooling, which are analyzed as case-study variants rather than main methods.
\end{itemize}
All configurations are managed through a hierarchical YAML setup so that only the intended stage-level
differences change across experiments.

\subsection{Data and Training Protocol}
The main data sources are AllNLI and STSBenchmark. AllNLI provides supervised pair relations for the
NLI warm-start and for constructing contrastive positives and hard negatives, while STSBenchmark
provides the graded similarity labels used in the final evaluation pipeline. For STSBenchmark, we
strictly separate \texttt{train}, \texttt{dev}, and \texttt{test} splits to avoid data leakage. When
an experiment requires pairwise contrastive or scored-pair supervision, these training pairs are
constructed automatically from the corresponding training resources before optimization.

Each configuration is run multiple times with different random seeds for robustness. The $i$-th seed
is defined as
\begin{equation}
    seed_i = seed_{start} + (i-1) \times seed_{step} .
\end{equation}
Each run starts from a clean initialization so that results from different methods remain directly
comparable.

\subsection{Evaluation}
After training, we evaluate the resulting checkpoints on the designated downstream benchmark and
aggregate the reported metrics across runs. For the STS-focused experiments in this report, the main
metric is cosine Spearman, following the original SBERT evaluation protocol
\cite{reimers-gurevych-2019-sentence-bert}.

\section{Methodology}
Our goal is to improve the original SBERT recipe without breaking compatibility with the 2019
\texttt{sentence-transformers} training stack. All reported best runs used the \texttt{bert-base-uncased} 
with mean pooling and maximum sequence length 128,
while the training curriculum is switched through \texttt{config.yaml}. 
The common backbone of all experiments is a three-stage view of sentence embedding learning. While the standard SBERT recipe relies on \textbf{Mean Pooling}:
\begin{equation}
\mathbf{h}(x)=\operatorname{MeanPool}(\operatorname{BERT}(x))
\end{equation}
we also investigated a structural variant using \textbf{Attention Pooling}. In this setting, the model learns to assign a weight $\alpha_i$ to each token embedding $\mathbf{h}_i$ via a trainable vector $w$:
\begin{equation}
\alpha_i = \frac{\exp(w^\top \mathbf{h}_i)}{\sum_j \exp(w^\top \mathbf{h}_j)}, \quad \mathbf{h}(x) = \sum_i \alpha_i \mathbf{h}_i
\end{equation}
This allows us to compare whether a learnable weighted aggregation of tokens can capture more distinctive semantics than a uniform average. 

\textbf{Regardless of the pooling method}, the training process follows a three-stage curriculum: we first warm-start the encoder on NLI, then optionally reshape the embedding geometry with an intermediate objective, and finally adapt to STS with direct similarity supervision.

\paragraph{Baseline SBERT curriculum.}
The baseline \texttt{sts\_nli} mode follows the original SBERT intuition
\cite{reimers-gurevych-2019-sentence-bert}: NLI classification teaches coarse semantic relations,
while STS fine-tuning calibrates the embedding space for graded similarity. In the NLI stage, the
model is trained with the original SBERT sentence-pair classifier, where the two sentence embeddings
$\mathbf{u}$ and $\mathbf{v}$ are combined as the concatenation
$[\mathbf{u};\mathbf{v};|\mathbf{u}-\mathbf{v}|]$ and passed to a softmax classification head. This
design is one of the key contributions of SBERT, because it transfers BERT into a siamese sentence
embedding framework while still exploiting strong supervised NLI signals. In the final STS stage, we
train on \texttt{sts-train.csv}, select by \texttt{sts-dev.csv}, and report only on
\texttt{sts-test.csv}; the objective is cosine regression on normalized STS labels:
\begin{equation}
\mathcal{L}_{\text{STS}}=\operatorname{MSE}\left(\cos(\mathbf{h}_i,\mathbf{h}_j), y_{ij}\right), \quad y_{ij}\in[0,1]
\end{equation}
This baseline is therefore not a trivial reference point, but a strong and carefully designed study
that remains the foundation for our later enhancements.

\paragraph{Enhancement I: AoE-style STS fine-tuning after NLI warm-start.}
Our first enhancement keeps the standard NLI warm-start but replaces the final STS regression loss
with an AoE-inspired ranking objective. The motivation comes from angle-optimized embeddings
\cite{li-li-2024-aoe}: cosine similarity is easy to optimize for moderately similar pairs, but its
gradient becomes weak in saturation regions close to $1$ or $-1$, which is undesirable for STS where
many sentence pairs differ only subtly. Instead of changing the encoder architecture, we keep the same
SBERT backbone and modify the STS supervision so that higher-scoring sentence pairs are ranked above
lower-scoring ones. In our implementation, this becomes a lightweight objective with two ranking terms,
one on cosine similarity and one on angle-derived similarity. For a pair $(x_i,x_j)$ with embeddings
$\mathbf{h}_i,\mathbf{h}_j$, the angle is
\begin{equation}
\Delta\theta_{ij}=
\operatorname{atan2}\left(
\begin{aligned}
&\left|\langle \mathbf{h}_i^{\Im}, \mathbf{h}_j^{\Re}\rangle
- \langle \mathbf{h}_i^{\Re}, \mathbf{h}_j^{\Im}\rangle\right|, \\
&\langle \mathbf{h}_i^{\Re}, \mathbf{h}_j^{\Re}\rangle
+ \langle \mathbf{h}_i^{\Im}, \mathbf{h}_j^{\Im}\rangle + \epsilon
\end{aligned}
\right)
\end{equation}
The model then applies a pairwise ranking loss so that sentence pairs with higher supervision scores
obtain smaller angles:
\begin{equation}
\mathcal{L}_{\text{angle}}=
\log \left(1+\sum_{y_p>y_q}\exp\left(\frac{\Delta\theta_p-\Delta\theta_q}{\tau_a}\right)\right)
\end{equation}
The real/imaginary split is only the implementation mechanism used to compute this angular relation;
the methodological focus is the ranking of pairwise angles rather than a ``complex-valued'' model in
itself.

The cosine term follows the same ranking idea on pairwise cosine scores, so the final objective is
\begin{equation}
\mathcal{L}_{\text{AoE-lite}}=
\lambda_{\text{cos}}\mathcal{L}_{\text{cos}}+
\lambda_{\text{angle}}\mathcal{L}_{\text{angle}} .
\end{equation}
In our implementation, STS labels are normalized to $[0,1]$ and used directly as ranking signals.
This branch therefore remains lightweight and fully compatible with the legacy 2019-style SBERT
training stack while still introducing geometry-aware similarity supervision. Unlike the heavier
hybrid branch discussed later, the reported \texttt{NLI+AoE+STS} result should be interpreted as an
AoE-style replacement for the STS loss rather than a separate contrastive curriculum built from merged
NLI and STS scored pairs.

For comparison, we also experimented with a stronger AoE-flavored hybrid that combines angle ranking
with an auxiliary contrastive term:
\begin{equation}
\mathcal{L}_{\text{AoE}}=\lambda_{\text{angle}}\mathcal{L}_{\text{angle}}+\lambda_{\text{cl}}\mathcal{L}_{\text{cl}}
\end{equation}
Intuitively, the angle loss improves local ranking sensitivity for fine-grained semantic similarity,
while the auxiliary contrastive term keeps positive pairs tightly aligned. However, as shown in the
case study section, this joint AoE+MNRL direction does not outperform the simpler AoE-lite branch in
our current setting.

\paragraph{Enhancement II: MNRL with UniReg before STS.}
One strong branch replaces the intermediate stage with hard-negative contrastive learning derived from NLI plus a
uniformity regularizer. The contrastive component is Multiple Negatives Ranking Loss (MNRL), which
can be viewed as a sentence-pair instantiation of the InfoNCE objective
\cite{rusak2025infonce}. In our setting, the matched positive is contrasted against all other
candidates in the batch together with hard negatives from NLI, making it closely related to the
supervised SimCSE formulation \cite{gao-etal-2021-simcse}. For each anchor sentence, its paired
positive is treated as the only correct target, while all other positives in the same batch and
explicit contradiction examples are used as negatives:
\begin{equation}
\mathcal{L}_{\text{MNRL}}=
-\log\frac{\exp(\operatorname{sim}(\mathbf{z}_i,\mathbf{z}_i^+)/\tau)}
{\sum_{k}\exp(\operatorname{sim}(\mathbf{z}_i,\mathbf{c}_k)/\tau)}
\end{equation}
Here $\mathbf{z}$ denotes L2-normalized sentence embeddings and $\mathbf{c}_k$ denotes the candidate
pool formed by the matched positive together with in-batch plus hard-negative examples. In our data
builder, NLI entailment pairs are used as positives and contradiction hypotheses for the same premise
are sampled as hard negatives, matching the supervised contrastive practice advocated by SimCSE.

While studying world models, we noticed the SIGReg regularizer used in LeWM \cite{maes2026lewm},
which prompted us to ask whether SBERT embeddings may also suffer from a related dimensional
collapse or cone effect; this heuristic observation led us to look for a lightweight geometric
regularizer and eventually adopt UniReg.
Following the alignment-uniformity perspective of contrastive representation learning
\cite{wang-isola-2020-alignment-uniformity}, we add a weak hyperspherical uniformity regularizer:
\begin{equation}
\mathcal{L}_{\text{uni}}=
\log\left(\frac{1}{|\mathcal{P}|}\sum_{u\neq v}\exp\left(-t\left\|\mathbf{z}_u-\mathbf{z}_v\right\|_2^2\right)+\epsilon\right),
\end{equation}
and optimize
\begin{equation}
\mathcal{L}_{\text{MNRL-UniReg}}=\mathcal{L}_{\text{MNRL}}+\lambda_{\text{uni}}\mathcal{L}_{\text{uni}}
\end{equation}
Minimizing this term encourages embeddings to spread more evenly on the unit hypersphere, which
reduces representation collapse and improves global geometry, while MNRL still preserves strong
positive-pair alignment. In our configuration, UniReg is intentionally lightweight
($\lambda_{\text{uni}}=10^{-3}$), so it acts as a global geometric prior rather than overriding the
task loss.
A small but consistent gain on geometry-level metrics was also observed on \texttt{sts\_dev}: the
effective rank increased from 196.6 to 204.7 out of an embedding dimension of 768 when UniReg was
added on top of MNRL.

\paragraph{Why these two enhancements are complementary.}
Although both branches improve cosine Spearman on STSBenchmark, they intervene at different levels of
the embedding geometry. AoE mainly changes the \emph{local ranking signal} by replacing saturated
cosine supervision with angle-based ordering on graded pairs. MNRL+UniReg instead keeps the familiar
contrastive objective but improves the \emph{global distribution} of embeddings through hard negatives
and hyperspherical spreading. This contrast makes the two methods a useful pair of lightweight
upgrades for a legacy SBERT codebase: one sharpens fine-grained similarity ordering, while the other
stabilizes the space in which retrieval-style similarity is computed.
A natural question is whether the two objectives can be combined. We did test a joint
\texttt{MNRL+AoE} stage, and it is analyzed in the case study section.

\section{Results and Analysis}
Table~\ref{tab:sts-benchmark-results} summarizes all completed STSBenchmark experiments under the cosine
Spearman metric, sorted by the improvements. Following Sentence-BERT
\cite{reimers-gurevych-2019-sentence-bert}, we use cosine Spearman as the primary metric because it
better reflects ranking quality for semantic similarity than Pearson correlation.

\begin{table}[ht]
\centering
\caption{Main STSBenchmark results for different methods.}
\label{tab:sts-benchmark-results}
\resizebox{\columnwidth}{!}{
\begin{tabular}{lcc}
\hline
Setting & Runs & Cosine Spearman (\%) \\
\hline
Baseline1: \texttt{STS} & 3 & $84.854 \pm 0.394$ \\
Baseline2: \texttt{NLI+STS} & 3 & $84.853 \pm 0.474$ \\
\texttt{NLI+AoE+STS} & 3 & $86.658 \pm 0.278$ \\
\texttt{NLI+MNRL+STS} & 3 & $86.749 \pm 0.127$ \\
\texttt{NLI+MNRL-UniReg+STS} & 3 & $\mathbf{86.792 \pm 0.126}$ \\
\hline
\end{tabular}
}
\end{table}

Besides the faithfully reproduced baselines, three observations are notable. 
First, all three strongest settings introduce an additional
geometry-shaping objective beyond the plain \texttt{STS} and \texttt{NLI+STS} baselines, and all of
them improve substantially over those two references. Second, the highest mean is obtained by
\texttt{NLI+MNRL+UniReg+STS}, whose cosine Spearman reaches $86.792\%$, marginally exceeding the rerun
\texttt{NLI+MNRL+STS} result of $86.749\%$ and the AoE branch at $86.658\%$. Because the gap over
plain rerun MNRL is small, we interpret UniReg as a promising lightweight improvement rather than a
decisive shift in training regime. Third, the AoE branch remains highly competitive, which suggests
that ranking-aware STS supervision is a strong alternative to contrastive-only geometry shaping in
this legacy SBERT setting. Notably, structural modifications such as Attention Pooling were also evaluated but showed diminishing returns compared to objective-level changes; these findings are detailed in the Case Study

Relative to the strongest plain baseline (\texttt{STS}, $84.854\%$), the gains are approximately
$+1.94$ points for UniReg, $+1.90$ for rerun MNRL, and $+1.80$ for AoE. Therefore, the current
evidence favors lightweight geometric regularization and ranking-aware objective design as the most
promising directions in this legacy SBERT setting.

\section{Case Study}

Table~\ref{tab:failure-sts-methods} summarizes three plausible variants that we tested but that still fall short of the best methods in this project. The common pattern is that additional complexity does not translate into better STS behavior unless the new objective is well aligned with the geometry that SBERT already learns.

\begin{table}[ht]
\centering
\caption{Performance of additional design variants on STSBenchmark.}
\label{tab:failure-sts-methods}
\resizebox{\columnwidth}{!}{
\begin{tabular}{lcc}
\hline
Setting & Runs & Cosine Spearman (\%) \\
\hline
Baseline: \texttt{NLI+STS} & 3 & $84.853 \pm 0.474$ \\
\texttt{NLI+MNRL+attn. pooling+STS} & 3 & $\mathbf{86.550 \pm 0.230}$ \\
\texttt{NLI+MNRL+AoE+STS} & 3 & $85.667 \pm 0.197$ \\
\texttt{NLI+MNRL+STS(CoSENT)} & 1 & $85.747 \pm 0.000$ \\
\hline
\end{tabular}
}
\end{table}

The \texttt{NLI+MNRL+attn. pooling+STS} variant presents an interesting case where structural complexity provides diminishing returns. While it achieves a significant gain over the plain \texttt{NLI+STS} baseline ($86.550\%$ vs. $84.853\%$), it fails to outperform the simpler \texttt{NLI+MNRL+STS} variant using mean pooling ($86.749\%$). This is consistent with the original SBERT finding that simple pooling over a strong siamese encoder is already a competitive inductive bias \cite{reimers-gurevych-2019-sentence-bert}. Under our limited data and compute budget, adding trainable parameters through an attention mechanism across the three-stage pipeline appears less useful than changing the training objective alone.

The most instructive failure is the joint \texttt{NLI+MNRL+AoE+STS} branch. Although both MNRL and AoE are individually effective, they emphasize different aspects of the embedding geometry. AoE is designed to improve fine-grained angular ordering in STS \cite{li-li-2024-aoe}, whereas contrastive objectives such as MNRL mainly reshape the space through batch-level alignment and uniformity \cite{gao-etal-2021-simcse}. When these signals are combined in a naive multi-objective stage, the optimization appears to become less coherent. Although the hybrid result is still better than the plain baselines, it is clearly worse than standalone \texttt{NLI+MNRL+STS} or \texttt{NLI+AoE+STS}. This suggests that local pairwise ranking and global contrastive shaping are not automatically complementary. Future work should therefore investigate how to combine their strengths more seamlessly.

The \texttt{NLI+MNRL+STS(CoSENT)} branch shows a milder limitation. Replacing the final regression loss with a ranking-style STS objective does not provide further improvement and instead pulls the MNRL branch slightly backward. This suggests that ranking supervision alone may be insufficient when the intermediate embedding geometry is not explicitly regularized.

Taken together, these failure cases reinforce the main lesson of the project: in a legacy SBERT stack, careful objective design contributes more than architectural complexity or naive objective composition.

\section{Conclusion}

This project revisited the original Sentence-BERT recipe under a legacy
\texttt{sentence-transformers} v0.2.4 codebase and a reproducible Slurm-based workflow. Rather than
targeting modern large-scale embedding leaderboards, our goal was to identify lightweight
objective-level improvements that remain compatible with a 2019-style SBERT training stack.

The final results show that training objective design matters more than pooling layer modification in this setting. 
Relative to the strongest plain baseline, all three enhanced branches improve cosine Spearman on STSBenchmark by around 1.8 to 1.9 points. 
Among them,
adding \texttt{MNRL+UniReg} stage achieves the highest mean score, 
while \texttt{AoE} remains highly competitive and demonstrates that ranking-aware STS supervision is an effective alternative to standard cosine regression. 
In contrast, naively combining MNRL and AoE in the same branch does not
improve further, suggesting that local angle-based ranking and contrastive geometry shaping may
interfere when jointly optimized.

Overall, the project shows that meaningful gains can still be obtained in a legacy SBERT codebase
through careful objective design, reproducible runtime configuration, and disciplined evaluation. A
natural next step is to further analyze why the hybrid \texttt{MNRL+AoE} setting underperforms,
improve the quality of negative sampling, and test whether simpler ranking objectives such as CoSENT
can provide a cleaner bridge between AoE-style supervision and contrastive learning.

\bigskip

\bibliography{aaai24}

\end{document}
